\newcommand{\figuraFamiliaPerfilesLosa}{%
\begin{figure}[H]
    \centering
    \begin{tikzpicture}[>=Latex,font=\small]
        \begin{axis}[
            width=11.8cm,
            height=7.4cm,
            xmin=0,xmax=4.5,
            ymin=0,ymax=1.08,
            axis lines=left,
            xlabel={$Z$},
            ylabel={$\rho(Z)/\rho(0)$},
            xtick={0,1,2,3,4},
            ytick={0,0.2,0.4,0.6,0.8,1},
            grid=major,
            grid style={gray!18},
            title={Transición entre autogravedad y gravedad externa},
            legend style={draw=none,fill=white,at={(0.98,0.98)},anchor=north east},
            samples=180]

            \addplot[very thick,blue!70!black,domain=0:4.5]
                {1/(cosh(x)^2)};
            \addlegendentry{$x=0$: losa autogravitante}

            \addplot[very thick,orange!85!black,dashed,domain=0:4.5]
                {1/(0.64*cosh(x/0.8+0.693147)^2)};
            \addlegendentry{$x=0.75$: régimen mixto}

            \addplot[very thick,red!75!black,dash dot,domain=0:4.5]
                {1/(0.2*cosh(x/0.447214+1.44364)^2)};
            \addlegendentry{$x=2$: gravedad externa dominante}

            \addplot[thick,green!45!black,densely dotted,domain=0:4.5]
                {exp(-4*x)};
            \addlegendentry{límite exponencial}

            \node[align=center,anchor=south east,text=black!70]
                at (axis cs:4.35,0.12)
                {mayor masa externa:\\perfil más concentrado};
        \end{axis}
    \end{tikzpicture}
    \caption{Familia de perfiles normalizados para una atmósfera sostenida
    por autogravedad y por la gravedad de una masa externa. Al aumentar el
    parámetro $x$, la gravedad planetaria domina y el perfil se aproxima a un
    decaimiento exponencial plano-paralelo.}
    \label{fig:familia-perfiles-losa}
\end{figure}%
}